
\documentclass{article}
\usepackage{colortbl}
\usepackage{makecell}
\usepackage{multirow}
\usepackage{supertabular}

\begin{document}

\newcounter{utterance}

\centering \large Interaction Transcript for game `eqbench', experiment `eq\_bench', episode 106 with reflexion{-}llama{-}3.1{-}8b{-}instruct.
\vspace{24pt}

{ \footnotesize  \setcounter{utterance}{1}
\setlength{\tabcolsep}{0pt}
\begin{supertabular}{c@{$\;$}|p{.15\linewidth}@{}p{.15\linewidth}p{.15\linewidth}p{.15\linewidth}p{.15\linewidth}p{.15\linewidth}}
    \# & \multicolumn{2}{c}{Player} && \multicolumn{2}{c}{Game Master} \\
    \hline

    \theutterance \stepcounter{utterance}  
    & & & \multicolumn{4}{p{0.6\linewidth}}{
        \cellcolor[rgb]{0.9,0.9,0.9}{
            \makecell[{{p{\linewidth}}}]{
                \texttt{\tiny{[P1$\langle$GM]}}
                \texttt{You have attempted this same task before and did not succeed. These are the plans you wrote after those attempts.} \\
\texttt{Trial \#1: New plan: I got stuck in a loop of analyzing the emotions of Terry and Cass without considering the context of the dialogue. I focused on the surface{-}level interactions between the two characters, but I didn't delve deeper into the underlying motivations and emotions that drive their actions. Next time, I will take a more nuanced approach by examining the power dynamics at play, the emotional triggers that are being pulled, and the underlying fears and insecurities that are being revealed. I will also pay closer attention to the language and tone used by each character, as well as the nonverbal cues and body language that can convey just as much information as the spoken words. By taking a more holistic approach, I hope to gain a deeper understanding of the emotional responses of Terry and Cass, and to arrive at a more accurate conclusion.} \\
\\ 
\texttt{Trial \#2: New plan: In my previous attempt, I focused on the surface{-}level emotions expressed by Terry and Cass, but I didn't consider the underlying motivations and power dynamics at play. I will take a more nuanced approach by examining the following:} \\
\\ 
\texttt{1. **Power dynamics**: I will analyze the relationship between Terry and Cass, including their roles, responsibilities, and areas of control. This will help me understand how their emotions are influenced by their positions and the expectations placed upon them.} \\
\texttt{2. **Emotional triggers**: I will identify the specific events, words, or actions that trigger strong emotions in Terry and Cass. This will enable me to understand what drives their reactions and how they respond to different stimuli.} \\
\texttt{3. **Underlying fears and insecurities**: I will explore the underlying fears, anxieties, and insecurities that are revealed through the dialogue. This will help me understand the deeper emotional currents that shape their interactions and responses.} \\
\texttt{4. **Language and tone**: I will pay closer attention to the language and tone used by each character, including their choice of words, tone, and pitch. This will help me understand their emotional states and how they convey their emotions through language.} \\
\texttt{5. **Nonverbal cues and body language**: I will consider the nonverbal cues and body language that are implicit in the dialogue, such as facial expressions, posture, and eye contact. This will provide additional insight into the emotional responses of Terry and Cass.} \\
\\ 
\texttt{By taking a more holistic approach, I hope to gain a deeper understanding of the emotional responses of Terry and Cass, and to arrive at a more accurate conclusion.} \\
\\ 
\\ 
\\ 
\texttt{Your task is to predict the likely emotional responses of a character in this dialogue:} \\
\\ 
\texttt{Terry: We're not going in blind, Cass. Not this time.} \\
\texttt{Cass: And who's gonna stop me? You?} \\
\texttt{Terry: If I have to. } \\
\texttt{Cass: Oh, right. The big, bad Terry Holloway. Always with a plan. } \\
\texttt{Terry: You're goddamn right. You know why? Because plans save lives.} \\
\texttt{Cass: Or maybe they just make you feel better about yourself. Like you're in control.} \\
\texttt{Terry: This isn't about control, Cass. It's about survival. } \\
\texttt{Cass: Survival? Or just your guilt?} \\
\texttt{Terry: That's enough.} \\
\texttt{Cass: Hit a nerve, did I? } \\
\texttt{Terry: You're good, Cass. But you're reckless. And that's a liability I can't afford.} \\
\texttt{Cass: And there it is. The threat. You gonna bench me, Terry?} \\
\texttt{Terry: If it means keeping you alive, yes.} \\
\texttt{Cass: I don't need your protection. I need to do my job.} \\
\texttt{{[}End dialogue{]}} \\
\\ 
\texttt{At the end of this dialogue, Terry would feel...} \\
\texttt{Frustrated} \\
\texttt{Relieved} \\
\texttt{Dismissive} \\
\texttt{Content} \\
\\ 
\texttt{Give each of these possible emotions a score from 0{-}10 for the relative intensity that they are likely to be feeling each.} \\
\\ 
\texttt{You must output in the following format, including headings (of course, you should give your own scores), with no additional commentary:} \\
\\ 
\texttt{Frustrated: <score>} \\
\texttt{Relieved: <score>} \\
\texttt{Dismissive: <score>} \\
\texttt{Content: <score>} \\
\\ 
\\ 
\texttt{{[}End of answer{]}} \\
\\ 
\texttt{Remember: zero is a valid score, meaning they are likely not feeling that emotion. You must score at least one emotion > 0.} \\
\\ 
\texttt{Your answer:} \\
            }
        }
    }
     \\ \\

    \theutterance \stepcounter{utterance}  
    & \multicolumn{4}{p{0.6\linewidth}}{
        \cellcolor[rgb]{0.561,0.737,0.561}{
            \makecell[{{p{\linewidth}}}]{
                \texttt{\tiny{[P1$\rangle$GM]}}
                \texttt{Frustrated: 8} \\
\texttt{Relieved: 0} \\
\texttt{Dismissive: 6} \\
\texttt{Content: 0} \\
            }
        }
    }
    & & \\ \\

    \theutterance \stepcounter{utterance}  
    & & & \multicolumn{2}{p{0.3\linewidth}}{
        \cellcolor[rgb]{0.9,0.9,0.9}{
            \makecell[{{p{\linewidth}}}]{
                \texttt{\tiny{[GM$|$GM]}}
                \texttt{Frustrated: 8} \\
\texttt{Relieved: 0} \\
\texttt{Dismissive: 6} \\
\texttt{Content: 0} \\
            }
        }
    }
    & & \\ \\

    \theutterance \stepcounter{utterance}  
    & & & \multicolumn{2}{p{0.3\linewidth}}{
        \cellcolor[rgb]{0.9,0.9,0.9}{
            \makecell[{{p{\linewidth}}}]{
                \texttt{\tiny{[GM$|$GM]}}
                \texttt{Frustrated: 8} \\
\texttt{Relieved: 0} \\
\texttt{Dismissive: 5} \\
\texttt{Content: 0} \\
            }
        }
    }
    & & \\ \\

    \theutterance \stepcounter{utterance}  
    & & & \multicolumn{2}{p{0.3\linewidth}}{
        \cellcolor[rgb]{0.9,0.9,0.9}{
            \makecell[{{p{\linewidth}}}]{
                \texttt{\tiny{[GM$|$GM]}}
                \texttt{game\_result = LOSE} \\
            }
        }
    }
    & & \\ \\

    \theutterance \stepcounter{utterance}  
    & & \multicolumn{4}{p{0.6\linewidth}}{
        \cellcolor[rgb]{0.95,0.95,0.95}{
            \makecell[{{p{\linewidth}}}]{
                \texttt{\tiny{[P1$\leftrightarrow$P]}}
                \texttt{Post{-}episode Reflexion (not sent to GM):} \\
\texttt{New plan: In my previous attempts, I focused on analyzing the surface{-}level emotions expressed by Terry and Cass, but I didn't consider the underlying motivations and power dynamics at play. I will take a more nuanced approach by examining the following:} \\
\\ 
\texttt{1. **Power dynamics**: I will analyze the relationship between Terry and Cass, including their roles, responsibilities, and areas of control. This will help me understand how their emotions are influenced by their positions and the expectations placed upon them.} \\
\texttt{2. **Emotional triggers**: I will identify the specific events, words, or actions that trigger strong emotions in Terry and Cass. This will enable me to understand what drives their reactions and how they respond to different stimuli.} \\
\texttt{3. **Underlying fears and insecurities**: I will explore the underlying fears, anxieties, and insecurities that are revealed through the dialogue. This will help me understand the deeper emotional currents that shape their interactions and responses.} \\
\texttt{4. **Language and tone**: I will pay closer attention to the language and tone used by each character, including their choice of words, tone, and pitch. This will help me understand their emotional states and how they convey their emotions through language.} \\
\texttt{5. **Nonverbal cues and body language**: I will consider the nonverbal cues and body language that are implicit in the dialogue, such as facial expressions, posture, and eye contact. This will provide additional insight into the emotional responses of Terry and Cass.} \\
\\ 
\texttt{However, I will also take a more specific approach by focusing on the following:} \\
\\ 
\texttt{* I will pay closer attention to the way Terry uses his authority and control to manage Cass's behavior, and how this affects Cass's emotional response.} \\
\texttt{* I will examine the way Cass uses her wit and sarcasm to push Terry's buttons and challenge his authority, and how this affects Terry's emotional response.} \\
\texttt{* I will consider the way the dialogue reveals Terry's guilt and Cass's recklessness, and how these underlying emotions shape their interactions and responses.} \\
\\ 
\texttt{By taking a more nuanced and specific approach, I hope to gain a deeper understanding of the emotional responses of Terry and Cass, and to arrive at a more accurate conclusion.} \\
            }
        }
    }
    & \\ \\

\end{supertabular}
}

\end{document}
